\documentstyle[%


righttag%


,11pt%


,amssymb%


,russian%               remove in Eng.


,izvru%                 Set izven% in Eng.


% ,izvru-ks             for brief communications


]{amsart}





%       Math definitions





\newcommand{\tts}[1]{}





%\hoffset=-15mm%                  For Eng.


\setcounter{page}{17}


\god{2000}


\nomer{3~(454)} %                        Remove in Eng.


%\nomer{Vol.\,44, No.\,3}%               Set in Eng.


%\str{\pageref{begin}--\pageref{end}}%  Set in Eng.


\udk{517.982:517.5}





\author{\it Г.В.\,Арутюнянц}


% NB:   ^^ remove in Eng.


\title{Об ограниченности некоторых


максимальных и мультипликаторных операторов}





\date{}


%\pagestyle{myheadings}%         Set in Eng.


\pagestyle{plain} %              Remove in Eng.


\begin{document}


%\thispagestyle{plain}%          Set in Eng.


\maketitle


\label{begin}





Пусть $\Theta$ --- произвольное подмножество отрезка $[0;\pi / 2]$.


Обозначим через $\Re_{\Theta}$ совокупность прямоугольников на плоскости


в направлениях из $\Theta$. Введенная совокупность является дифференциальным


базисом.





Рассмотрим теперь общую проблему о том, {\it как связаны область значений


$p \in [ 1;+\infty]$, для которых базис $\Re_{\Theta}$


дифференцирует $L^p (R^2)$, с конкретным видом множества $\Theta$}.


При этом под дифференцированием класса будем понимать дифференцирование


интегралов всех функций из данного класса.





Классическая теорема Б.\,Йессена, Ю.\,Марцинкевича и А.\,Зигмунда [1] влечет


дифференцируемость $L\log^{+} L$, если $\Theta$ --- конечное множество.


Позже А.\,Зигмунд заметил, что если $\Theta= [ 0;\pi / 2]$, то


$\Re_{\Theta}$ не является даже плотностным, т.\,е. $\Re_{\Theta}$


не дифференцирует $L^{\infty}$. В рамках тех же рассуждений можно получить


и более сильный результат. {\it


Пусть $\Theta \subset [0;\pi/2]$ и замыкание этого множества имеет


положительную меру. Тогда $\Re_{\Theta}$ также является неплотностным}


([2], с.\,228). Таким образом, интерес представляет случай, когда замыкание


множества $\Theta$ имеет нулевую меру.





Пример бесконечного множества $\Theta$ с хорошими свойствами $\Re_{\Theta}$


впервые был предложен в [3], где рассмотрен случай, когда


$\Theta$ является лакунарной последовательностью, например,


$\Theta=\{\pi 2^{-k} \}_{k=1}^\infty$. Используя оригинальную


геометрическую идею, Й.\,Стремберг доказал, что $\Re_{\Theta}$


дифференцирует $L^2 (\log^{+} L) ^{4+\varepsilon}$,


$\varepsilon > 0$. Позже в [4] было показано, что


$\Re_{\Theta}$ дифференцирует $L^2$. Наконец в


[5] было доказано, что $\Re_{\Theta}$ дифференцирует $L^p$, $p>1$.


Предложенный в [5] подход основан на современных достижениях анализа


Фурье таких, как интерполяция в пространствах со смешанной нормой


(так называемый ``bootstrap argument''), рандомизация, теорема Марцинкевича о


мультипликаторах Фурье и пр. Вопрос о геометрическом доказательстве теоремы


А.\,Нагеля, И.\,Стейна и С.\,Вейенгера [5]


является важной открытой проблемой.


Ее решение позволило бы найти точный класс дифференцирования. Пока лишь


известно, что этот класс не шире, чем $L \log^{2} L$.





Результаты из [5] были обобщены в [6], где


рассматривались множества, получаемые путем итераций лакунарных направлений.


Было показано, что в случае конечных итераций соответствующий базис


$\Re_{\Theta}$ также дифференцирует $L^p$, $p>1$. В случае бесконечных


итераций соответствующее множество $\Theta$ приобретает вид Канторова


множества, которое можно записать в форме


$$


\Theta = \bigg\{\sum_{i=1}^{\infty} \varepsilon_i \omega_i,


\ \; \varepsilon_i\in\{0,1\};


\ \; \omega_n>0;\ \; \sum_{i=n+1}^{\infty} \omega_i < \omega_n,


\ \; n \in N;\ \;


\sum_{i=1}^{\infty} \omega_i = \pi / 2 \bigg\}.


$$


В [6] доказано, что если


$\omega_{n+1} / \omega_n=q$, $0<q<1/2$, $n \in N$, то $\Re_{\Theta}$


не дифференцирует $L^p$, где


$$


p < 1+\log 2 / \log q^{-1}.


$$


Классическое Канторово множество на $[0;\pi/2]$ получается, если


$\omega_{n+1} / \omega_n=1/3$, $n \in N$.





В данной статье изучается предельный случай


$\lim\limits_{n \to \infty} \omega_{n+1} / \omega_n=1/2$.


Показывается (теорема 1), что соответствующий базис не является плотностным.





Определим максимальный оператор $ {\cal M}_{\Theta}f $ по правилу


$$


({\cal M}_{\Theta}\, f) (x) = \sup_{x \in I \subset


\Re_{\Theta}}| I |^{-1} \int_I | f(y) | \, dy,


$$


где $f \in L_{\text{loc}} (R^2)$.





В [7] доказано, что в случае классического Канторова множества


соответствующий дифференциальный базис $\Re_{\Theta}$ не дифференцирует


$L^p$, $1 \le p \le 2$.





Такое альтернативное поведение дифференциальных свойств базисов


$\Re_{\Theta}$ позволяет высказать предположение:


{\it


либо базис $\Re_{\Theta}$ дифференцирует $L^2$, либо не является


плотностным.}





Доказательство теоремы 1 основано на изучении арифметической структуры


множеств Канторового типа. А именно, выбирается некоторая регулярная


последовательность направлений, для которой используется конструкция


Безиковича. Несмотря на то, что используются относительно элементарные


рассуждения, полученный результат может быть применен в достаточно


нетривиальных вопросах, например, в теории мультипликаторных операторов


Фурье (теорема 2). Отметим, что вопрос об арифметической природе эффектов,


наблюдаемых в теории мультипликаторов, достаточно часто обсуждается в


различных работах. В частности, интересные рассуждения можно найти в


([8] и [9], с.\,162).





\begin{thm}


Если


\begin{equation}


\lim_{n \to \infty} \frac {\omega_{n+1}} {\omega_n} =


\frac {1}{2},


\end{equation}


то $\Re_{\Theta}$ не является плотностным.


\end{thm}





Пусть $\wh R$ --- произвольный прямоугольник на плоскости. Продлим пару его


сторон в обе стороны, а также проведем два отрезка параллельных паре других


сторон прямоугольника $\wh R$, таким образом, чтобы получилось три


прямоугольника одинаковой меры. Обозначим через $R$ средний прямоугольник,


а через $\wt R$ --- объединение двух оставшихся. В этих терминах


сформулируем основную лемму, из которой теорема 1 получается стандартными


рассуждениями (см. [2], с.\,370).





\begin{lem}


Пусть $\Theta=\{\Theta_n\}_{n=1}^\infty,$ $\Theta_n \downarrow 0$ и


\begin{equation}


\lim_{n \to \infty} \frac {\Theta_{n+1}} {\Theta_n} =1.


\end{equation}


Тогда для некоторого $q>0$ и любого $ \eta >0$ можно найти множество $E$ и


прямоугольники $\{R_k \}_{k=1}^m$ со следующими свойствами$:$


\begin{itemize}


\item[(i)] направление $R_s \in \Theta$, $s=1,\dotsc,m$,


\item[(ii)] $| E | \le \eta \sum\limits_{k=1}^m | R_k |$,


\item[(iii)] $| \wt R_k \cap E | \ge q |


\wt R_k |$, $| R_i \cap R_j | = \emptyset$, $i \ne j$.


\end{itemize}


\end{lem}





Из последних трех свойств следует, что для некоторого $q>0$ можно найти


последовательность множеств $\{E_n \}_{n=1}^{\infty}$, для которой


\begin{equation}


\lim_{n \to \infty} \frac {| \{({\cal M}_{\Theta}


\chi_{E_n}) (x) > q \}


|} {| E_n |} =\infty.


\end{equation}





Используя (3) и критерий Буземана--Феллера [10], устанавливаем, что базис


$\Re_{\Theta}$ не является плотностным.





Лемма 1 является модификацией леммы Ч.\,Феффермана из [11]. Она была получена


А.M.\,Сто\-ко\-ло\-сом [12] при условии, что


$$


\Theta_n\downarrow 0,\quad


\limsup_{n \to \infty} \Theta_{n+1} / \Theta_n=1


$$


и $\{\Theta_n^{-1} \}_{n=1}^{\infty}$ является выпуклой. Условия


леммы 1 позволяют доказать ее, используя те же рассуждения.





Таким образом, необходимо доказать, что при условии (1) можно выбрать


последовательность точек из $\Theta$, обладающих свойством (2).





\begin{lem}


Пусть $\lim\limits_{n \to \infty} \omega_{n+1} / \omega_n =1 / 2$.


Тогда существует последовательность точек $\{\Theta_n\}$ из $\Theta$


такая, что $\Theta_n \downarrow 0$ и


$$


\lim_{n \to \infty} \frac {\Theta_{n+1}} {\Theta_n} =1.


$$


\end{lem}





{\bf Доказательство.} Введем обозначения


$$


\alpha_n= \omega_{n+1} / \omega_n;\ \;


\wt {\alpha} _n =\inf (\alpha_i, i \ge n);\ \; K_n= \Theta_n \cap


[ \gamma_{n+1}; \gamma_n ],\ \; n \in N,


$$


где $\gamma_n \downarrow 0$, $\lim \gamma_{n+1} /


\gamma_n =1$, $\gamma_n \notin \Theta$. Заметим, что такая последовательность


$\{\gamma_n \}_{n=1}^{\infty}$ существует, т.\,к. $\Theta$ нигде не плотно.


Не ограничивая общности, можно считать $K_1 \ne \emptyset$. Пусть


$\Theta_1 = \inf K_1 $ и $\Theta_2,\dotsc,\Theta_n$ уже выбраны. И пусть


$\Theta_n$ принадлежит некоторому $K_s$, $s \in N$.


Выберем $\Theta_{n+1}$ следующим образом:





(a) eсли $\Theta_n =\sup K_s$, то $\Theta_{n+1} =\inf K_s$;





(b) если $\Theta_n =\inf K_s$, то $\Theta_{n+1} =\sup


\cup_{i=s+1}^{\infty} K_i$.





Очевидно, $\Theta_n \downarrow 0$, $\Theta_n \in \Theta$.


Далее рассмотрим два случая.





1) $\Theta_n$ и $ \Theta_{n+1}$ принадлежaт одному множеству, например,


$K_p$ или двум соседним $K_p$, $K_{p+1}$. Тогда


$$


1 \ge \frac {\Theta_{n+1}}{\Theta_n} > \frac {\gamma_{p+1}}{\gamma_p},


$$


и (2) следует из этой оценки.





2) В противном случае, используя геометрическую структуру множества


$\Theta$,


отметим, что $\Theta_n$ и $\Theta_{n+1}$ являются концами интервала $I$,


принадлежащего дополнению множества $\Theta$. Тогда существуют


$q=q (n)$, $n \in N$, и $\eta \ge 0$ такие, что


$$


\Theta_n = \eta+\omega_q; \quad \Theta_{n+1} = \eta+


\sum_{i=q+1}^{\infty} \omega_i.


$$


Следовательно,


$$


1 \ge \frac {\Theta_{n+1}} {\Theta_n}


\ge \frac {\sum\limits_{i=q+1}^{\infty} \omega_i} {


\omega_q} =


\sum_{i=q+1}^{\infty} \prod_{j=q} ^{i-1} \alpha_j =


\sum_{i=q+1}^{\infty}\wt {\alpha}_q^{i-q} =


\sum_{s = 1}^{\infty}\wt {\alpha}_q^s= \frac {\wt \alpha_q}


{1-\wt \alpha_q}.


$$





Осталось заметить, что $q=q(n) \to \infty$ при $ n \to \infty$.


Лемма 2 и теорема 1 доказаны.


\medskip





{\bf Замечание.}


Если $\omega_n=2^{-n}\varepsilon_n$, где $\varepsilon_n


\downarrow 0$,


$\lim\limits_{n \to \infty}\varepsilon_{n+1} / \varepsilon_n=1$,


то по теореме 1 соответствующий базис $\Re_{\Theta}$ является неплотностным,


однако нетрудно показать, что $\text{mes\,} \Theta =0$.





Укажем применение полученного результата в теории мультипликаторов.


Определим в терминах преобразования Фурье мультипликаторный


оператор $\widehat{T_E f(x)} = \chi_E (x)


\wh {f} (x)$, $f \in L^2 \cap L^p (R^2)$, $p>1$,


где $E$ --- измеримое множество. Говорят, что множество $E$ является


$L^p$-муль\-ти\-пли\-ка\-то\-ром,


если соответствующий мультипликаторный оператор ограничен в $L^p$.





Хорошо известно, что многоугольник с конечным числом сторон является


$L^p$-муль\-ти\-пли\-ка\-то\-ром для $1<p< \infty$.


С другой стороны, в


[11] доказано, что единичный круг является мультипликатором только в


тривиальном случае $p=2$. Поэтому возник интерес к рассмотрению множеств,


промежуточных в некотором смысле между обычным многоугольником и кругом


([2], с.\,369; [12]).





Пусть $\Theta^*$ --- периодическое продолжение $\Theta $ на $[0;2 \pi]$,


а $\Phi=\{(\cos \phi, \sin \phi) : \phi \in \Theta^*


\}$. Пусть $P_{\Theta}$ образован пересечением полуплоскостей, содержащих


единичный шар, границы которых касаются единичной окружности в точках $\Phi$.


Очевидно, $P_{\Theta}$ является выпуклым многоугольником, стороны


которого имеют нормали в направлении множества Канторового типа. Отсюда


стандартной техникой ([2], с.\,362; [11]) получается





\begin{thm}


Если


$$


\lim_{n \to \infty} \frac {\omega_{n+1}} {\omega_n} =


\frac {1} {2},


$$


то $T_{P_{\Theta}}$ ограничен только в $L^2 (R^2)$.


\end{thm}





В заключениe автор выражает благодарность А.М.\,Стоколосу за плодотворные


обсуждения и полезные советы, а также П.\,Шегрену за интерес, проявленный


к работе.





\begin{thebibliography}{99}





\bibitem{1}


Jessen B., Marcinkiewicz J., Zygmund A. {\it Note on the


differentiability of multiple integrals} // Fund. Math. -- 1935.


-- V.\,25. -- P.\,217--234.


\bibitem{2}


De~Guzman M. {\it Real variable methods in Fourier analysis}.


-- Amsterdam -- N.~Y.--Oxford: North-Holland Publ. -- 1981. -- 392 p.


\bibitem{3}


Str\"omberg J.O. {\it Weak estimates on maximal functions


with rectangles in certain directions} // Ark. Math. -- 1977.


-- V.\,15. -- P.\,229--240.


\bibitem{4}


Cordoba A., Fefferman R. {\it On differentiation of integrals} //


Proc. Nat. Acad. Sci. USA. -- 1977. -- V.\,74. -- \No\,6. -- P.\,2211--2213.


\bibitem{5}


Nagel A., Stein E.M., Wainger S.


{\it Differentiation in lacunary directions} //


Proc. Nat. Acad. Sci. USA. -- 1978. -- V.\,75. -- P.\,1060--1062.


\bibitem{6}


Sj\"ogren P., Sj\"olin P. {\it Littlewood--Paley decompositions and Fourier


multipliers with singularities on certain sets} //  Ann. Inst. Fourier.


-- 1981. -- V.\,31. -- \No\,1. -- P.\,157--175.


\bibitem{7}


Katz N.H. {\it Counterexamples for maximal operators over Cantor sets


of directions} // Math. Research Letters. -- 1996. -- V.\,3.


-- \No\,4. -- P.\,527--536.


\bibitem{8}


Hare K.E., Klemes I. {\it Properties of Littlewood--Paley sets} //


Math. Proc. Camb. Phil. Soc. -- 1989. -- V.\,83. -- P.\,485--495.


\bibitem{9}


Edwards R.E., Gaudry G.I.


{\it Littlewood--Paley and multiplier theory}. -- Berlin:


Springer, 1977. -- 212~p.


\bibitem{10}


Buseman H, Feller W.


{\it Zur Differentiation das Lebesgueschen Integrale} //


Fund. Math. -- 1934. -- Bd.\,22. -- S.\,226--256.


\bibitem{11}


Fefferman C. {\it The multiplier problem for the ball} // Ann.


Math. -- 1971. -- V.\,94. -- \No\,2. -- P.\,330--336.


\bibitem{12}


Стоколос A.M. {\it К вопросу М.\,Гусмана о мультипликаторах Фурье


полигональных областей} // Сиб. матем. журн. -- 1995. --


T.\,36. -- \No\,6. -- C.\,1392--1398.





\end{thebibliography}





\vskip 2cm





\postup{\it Институт математики, экономики}{$\qquad$\it Поступили}


\postup{\it и механики $($г.\,Одесса, Украина$)$}{{\it первый вариант} 22.04.1997}


\postup{}{{\it окончательный вариант} 22.04.1999}


\label{end}


\end{document}


